% IEEEtran journal layout: two columns. Standalone course-review source.
% Template: https://www.michaelshell.org/tex/ieeetran/
% Bibliography overlap with teammates remains to be checked before submission.
\documentclass[journal]{IEEEtran}
\usepackage[T1]{fontenc}
\usepackage{amsmath,amssymb,booktabs,tabularx}
\usepackage{hyperref}
\hypersetup{colorlinks=true,allcolors=black}
\urlstyle{same}
\pagestyle{plain}
\begin{document}
\title{Learning Social Influence: A Review of Hyperbolic Representations, Seed Selection, and Reproducible Evaluation}
\author{Liu Shengqi\\Individual Related-Work Review, October 2026}
\maketitle
\begin{abstract}
Influence maximization connects social-network structure, diffusion dynamics, and the selection of a limited set of initial users. This review examines the transition from model-based optimization and local discounting to learned influence representations. It compares classical baselines, the DeepIM family, sheaf neural diffusion, graph neural bandits, and hyperbolic learning. HIM provides the focus for a course reproduction, while related diffusion-prediction research clarifies the opportunities and limits of geometric modeling. Preliminary experiments illustrate why model selection, independent evaluation, and reproducibility must be treated separately. The central argument is that a useful representation must support the quality of an entire seed set, rather than merely identify individually prominent users.
\end{abstract}

\section{Problem and Review Perspective}
Influence maximization asks which small group of users should initially receive information if the objective is broad subsequent diffusion. For a graph $G=(V,E)$ and a seed budget $k$, the usual formulation is
\begin{equation}
S^*\in\arg\max_{S\subseteq V,\ |S|\leq k}\sigma(S),
\end{equation}
where $\sigma(S)$ denotes expected activation under a declared diffusion process. This objective concerns a set: two individually influential users may reach almost the same audience. The problem therefore differs from sorting users by degree, popularity, or a learned node score.

My project contribution is to reproduce HIM and examine its representation and seed-selection choices. This review compares what approaches observe, how they handle redundant reach, and what evidence supports their claims. Twelve publications are discussed, including six published in 2024--2026. Their roles range from foundational optimization to recent learning and temporal diffusion research; they do not all need to be implemented in the course project.

\section{Optimization, Locality, and Redundant Reach}
Kempe, Kleinberg, and Tardos established an influential formulation under stochastic diffusion models \cite{kempe}. For independent cascade and the standard linear threshold setting, expected spread is monotone and submodular. With an exact influence oracle and a cardinality constraint, greedy selection obtains the $1-1/e$ approximation guarantee. These assumptions matter: the guarantee is not automatically inherited by an arbitrary learned predictor or a finite, noisy simulation procedure.

CELF exploits diminishing returns to reduce unnecessary marginal-gain evaluations \cite{celf}. Its main contribution is computational reuse. Under an identical deterministic objective and tie-breaking rule, lazy evaluation implements greedy selection rather than a different quality objective. Matching Greedy and CELF seed sequences can therefore be an implementation check. Independent Monte Carlo calls for every candidate can obscure that check because changing estimates change apparent gain rankings. Separating oracle construction from selection time also avoids hiding simulation costs.

Degree discount updates local scores after a seed is selected \cite{discount}. In the common small, uniform-probability setting, selected neighbors reduce the estimated usefulness of a candidate. This cheaply addresses overlap without estimating every global marginal gain. Its limitation is the gap between local exposure and multi-hop reach: a bridge between regions may be important despite moderate degree. The lesson for learning methods is that an overlap adjustment can be valuable even after a strong ranking has been obtained. Historical age does not remove the usefulness of degree discount as a clearly specified reference baseline.

\section{Learning Influence and Selecting Sets}
DeepIM learns a latent representation of seed sets together with a data-driven diffusion model \cite{deepim23}. It links representation learning to combinatorial selection rather than treating node classification as a substitute for influence maximization. Its budget formulation also accommodates node-centrality constraints. The related 2024 Neural Networks publication emphasizes accelerated inference and develops the DeepIM framework further \cite{deepim24}. These publications belong to the same method family; citing both does not constitute two independent confirmations of its performance.

The DeepIM direction raises a practical question: when is learning worth its initial cost? Fast inference can be attractive for repeated decisions, but one graph and one seed budget may not amortize training and supervision generation. My interpretation is that comparisons should report the number of intended deployments and separate training, inference, and evaluation costs. Another question is whether training examples cover the combinations encountered during optimization. Prediction accuracy on familiar seed sets alone does not establish reliable ranking of unfamiliar candidate sets.

DeepSN uses sheaf neural diffusion to model diverse influence patterns and introduces optimization that accounts for overlapping influence \cite{deepsn}. Its significance is the connection between a richer propagation representation and set selection. Compared with a simple ranking, it places more modeling capacity on interactions. However, extra capacity does not itself establish better generalization. A controlled study must align observed history, training graphs, budgets, and evaluation information before attributing a performance difference to representation.

Graph neural bandits address another information regime. Feng, Tan, and Cautis study multi-round campaigns with limited knowledge of network topology \cite{bandit}. Their IM-GNB framework combines exploration and exploitation when choosing influencers and aims to reach distinct users across rounds. It distinguishes learning from fixed historical samples from actively acquiring feedback. Its campaign objective differs from single-shot static selection. An offline HIM experiment cannot establish superiority over a bandit policy without specifying rounds, feedback, and repeated exposure.

\section{Why Hyperbolic Geometry Is Relevant}
Poincar\'e embeddings represent hierarchical relationships in negatively curved space \cite{poincare}. Their design captures aspects of similarity and hierarchy and uses Riemannian optimization to learn compact representations. Nickel and Kiela subsequently studied hierarchy learning in the Lorentz model, reporting computational and numerical advantages over their Poincar\'e formulation \cite{lorentz}. These works establish geometric tools for hierarchy learning; they do not themselves prove that hyperbolic embeddings optimize social influence.

For social diffusion, the appeal is an inductive bias: users may occupy different influence levels, while nearby users may share audiences. A radial coordinate can support hierarchical ordering and angular position can encode relatedness. Nevertheless, this interpretation is a modeling hypothesis, not a direct measurement of causal persuasion. Dense communities, repeated participation, and sampling bias could also shape the geometry. Evaluation should ask whether geometry improves set outcomes beyond degree-informed initialization and whether any advantage persists outside training conditions.

HIM learns hyperbolic influence representations from social edges and observed propagation activations, then selects seeds using distance information \cite{him}. Its adaptive sliding window (ASW) updates candidate scores to reduce redundant choices. The MD variant selects users with the smallest origin distances. This makes representation learning and overlap-aware selection separately examinable. The model does not require diffusion parameters as direct learning inputs; simulated histories can nevertheless depend on a particular generating mechanism. Diffusion-model-agnostic design should not be read as guaranteed robustness to arbitrary distribution shift.

More precisely, HIM uses a Lorentz model with relation-specific rotations for social edges and propagation edges. Negative-sampling objectives bring observed related pairs together while separating sampled negatives. An additional activation regularizer moves users that successfully activate others toward the origin. If $\mathbf{x}_u$ is user $u$'s learned representation and $\mathbf{o}_{\mathcal L}$ is the Lorentz origin, the influence-strength proxy is
\begin{equation}
\mathrm{LDO}_u=d_{\mathcal L}^{2}(\mathbf{x}_u,\mathbf{o}_{\mathcal L}).
\end{equation}
Smaller LDO is interpreted as stronger influence. This is a learned proxy, not a causal estimate of how many real people a user persuades. The structural and propagation objectives provide complementary evidence: the former describes who is connected, while the latter describes who activated whom in the observed cascades.

ASW addresses a different part of the problem from LDO. It starts with the low-LDO ranking, keeps a window of approximately $w=\beta k$ promising candidates, and updates neighboring candidates after each seed is chosen. The penalty is larger for candidates that are close to the selected seed in hyperbolic distance, so the method tries to avoid selecting users whose audiences may overlap. The search is restricted to graph neighbors in the current candidate window, making ASW a lightweight heuristic rather than an exact marginal-gain calculation. It has no general guarantee of finding the optimal seed set. Comparing HIM with HIM-MD therefore tests the added value of this overlap-aware heuristic while keeping the learned representation fixed.

The paper evaluates on Cora-ML (2,810 nodes), Power Grid (4,941), Facebook (5,908), GitHub (37,700), and YouTube (1,134,890), with 1\%, 5\%, and 10\% seed budgets. It uses IC/LT settings with specified parameters and WIC/WLT settings whose parameters are withheld from the methods. For each setting, 30 simulated propagation instances provide learning data; average spread is evaluated over 100 rounds. The unknown-parameter setting is still a controlled simulation: it does not mean the paper observed real opinions or real-world cascades. In particular, the result depends on the chosen generator, fixed parameter instance, training cascades, and evaluation protocol.

There is also a reproduction-sensitive paper/code distinction. The paper describes WIC edge probabilities sampled from $[0,0.2]$, while the bundled Cora-ML probability file used in this project spans approximately $[0.1,0.2]$. The paper says WLT parameters are sampled once and held fixed across propagation realizations; the repository's WLT simulator resamples node thresholds during each simulation. Thus, results from the bundled code should be labeled with the actual files and runtime behavior, rather than presented as an exact reproduction of every manuscript setting.

THGNets offers a nearby but different use of hyperbolic learning \cite{thgnets}. It predicts the next activated user using temporal hypergraphs, social-network structure, and temporal consistency constraints. Its relevance is the treatment of hierarchical and higher-order interactions in diffusion data. Prediction accuracy is not a seed-set optimization guarantee: a model can predict frequent continuations well without identifying complementary initial users. Transferring such a predictor to maximization requires a selection procedure and evaluation of complete cascades.

\begin{table}[t]
\caption{Different learning settings require different comparisons.}
\centering\footnotesize
\begin{tabularx}{\columnwidth}{@{}p{0.22\columnwidth}XX@{}}
\toprule
Method family & Main information & Evaluation issue \\
\midrule
Greedy/CELF & Specified diffusion oracle & Marginal gains and oracle cost \\
Degree discount & Local graph statistics & Limited multi-hop information \\
DeepIM/DeepSN & Training data and graph & Set interactions and training cost \\
IM-GNB & Sequential feedback & Exploration across campaigns \\
HIM & Graph and propagation history & Ranking plus overlap \\
THGNets & Temporal cascades & Prediction, not seed optimization \\
\bottomrule
\end{tabularx}
\end{table}

\section{Temporal Structure and Evaluation Limits}
Yanchenko, Murata, and Holme review temporal influence maximization, including seed activation timing and evolving contacts \cite{temporal}. This matters if our dataset is constructed from timestamped communications. Static aggregation may create a path whose contacts occur in an impossible temporal order. A large static simulated cascade does not demonstrate that the same diffusion could occur within the observed window. A temporal extension must specify activation timing, available past information, and future network uncertainty.

Fair comparison also requires an information boundary. A classical baseline accessing true edge probabilities receives information unavailable to a history-only learner. It can remain useful as a known-parameter reference, but should not be described as an equally informed competitor. Conversely, withholding parameters does not excuse using final evaluation outcomes for repeated tuning. Shared generating parameters, separate evaluation randomness, and identified validation decisions make comparisons interpretable.

\section{Preliminary Reproduction as Methodological Evidence}
I ran the author implementation on Cora-ML with 2,810 nodes and 7,981 undirected edges, using 30 simulated propagation instances and 281 seeds (10\%). One WIC+ASW run used training seed 42 and stopped at epoch 54. Its highest validation coverage was 38.15\% at epoch 43. Independent evaluation of the saved last-epoch seed set over 1,000 simulations produced mean spread 1,057.107, or 37.6195\% coverage. Reloading the saved model recovered identical seeds; repeating evaluation with seed 1042 reproduced the stored mean and standard deviation exactly.

This single-run evidence supports successful execution and artifact consistency; by itself, it does not measure variation across training runs, establish an ASW advantage, or reproduce every paper table. The highest validation result and independent evaluation refer to different checkpoints: the saved model is from the final executed epoch, not the epoch with the best validation score. Future comparisons should predefine checkpoint selection and preserve the corresponding model and seeds. Scripts reporting only a maximum during training leave this distinction unresolved.

My environment uses macOS CPU execution and compatible dependencies. The diffusion constructor was adjusted to preserve the caller's random stream when no explicit seed is supplied. Caches were checked to avoid reusing a smaller smoke-test sample. These choices support repeatability, but runtime cannot be directly compared with GPU measurements. Observed spread is conditional on the supplied simulation parameters, not evidence that real public opinions changed.

\subsection{Paired Selection Ablation}
To isolate selection from representation learning, I loaded the same saved Cora-ML checkpoint and compared HIM-ASW with HIM-MD using the same 281-seed budget. Both methods selected 281 unique users. For each of 1,000 evaluation replicates, Python and NumPy were initialized with the same replicate-specific seed for both methods; the evaluation used the same fixed WIC probability file. ASW achieved mean spread 1,057.466 ($37.6322\%$ coverage), versus 1,054.060 ($37.5110\%$) for MD. The paired ASW--MD difference was $+3.406$ activated users, with a 95\% paired $t$ interval of $[1.943,4.869]$, equivalent to $+0.1212$ percentage points of coverage (interval $[0.0692,0.1733]$). ASW was higher in 540 replicates, tied in 20, and lower in 440.

This is evidence of a small conditional advantage for ASW in this particular checkpoint and simulator, not evidence of a practically large effect or of robustness across training seeds. The paired interval captures variation over the 1,000 stochastic evaluations only; it does not capture training-run uncertainty. Resetting both pseudorandom generators to the same replicate seed aligns their initial random streams, but different seed sets can cause the simulator to consume random numbers along different paths. This is therefore a matched-seed evaluation, not a claim that every edge attempt used identical random variables. The seed sets, hashes, software versions, evaluation protocol, and aggregate statistics are stored in \texttt{results/HIM/paired\_selection\_cora\_ml\_wic\_2026-10-11.json}.

\subsection{Ciao Dataset Adaptation and Initial Runs}
I adapted the Ciao trust graph to HIM's input format and generated a fixed symmetric weighted-cascade probability bank, with $p(u,v)=1/\max(\deg(u),\deg(v))$ on each undirected edge. This is a declared synthetic process on a static trust topology, not an estimate of real Ciao transmission probabilities or observed opinion adoption. One smoke run trained for one epoch from only two sampled propagation instances, selected 366 users (the rounded 5\% budget), and successfully saved a checkpoint and seed set. The graph contains 7,317 users and 85,205 undirected edges. The smoke-training validation rounds are not independent test results.

I then reloaded that checkpoint and evaluated its ASW and MD seed sets on the same fixed probability bank for 1,000 matched-seed Monte Carlo replicates, using evaluation seed base 1042. ASW reached mean spread 714.156 users (9.7602\% coverage), and MD reached 713.378 (9.7496\%). The paired difference was $+0.778$ activated users, with a 95\% paired $t$ interval of $[-0.289,1.845]$; ASW was higher in 517 replicates, tied in 17, and lower in 466. Since the interval includes zero and the checkpoint was trained on only two instances for one epoch, this is strictly pipeline evidence, not a performance claim. The full seed lists, per-replicate outcomes, hashes, and settings are stored in \texttt{results/HIM/paired\_selection\_ciao\_wic\_smoke\_2026-10-11.json}.

After replacing a missing Python-environment symlink with a dedicated Python 3.12.14 environment and reinstalling the pinned macOS dependencies, I trained a larger first model with 30 propagation instances, training seed 42, and a 5\% budget. Training allowed up to 200 epochs with 100 online spread-evaluation rounds per epoch and patience 10; it stopped after epoch 39 and selected the epoch-28 checkpoint. For this checkpoint, ASW achieved mean spread 733.197 users (10.0205\% coverage), and MD achieved 732.357 (10.0090\%) over 1,000 independent matched-seed evaluation replicates. The paired difference was $+0.840$ users, with a 95\% paired $t$ interval of $[-0.843,2.523]$; ASW was higher in 504 replicates, tied in 20, and lower in 476. This per-checkpoint interval includes zero; the across-training-seed analysis below is the appropriate initial estimate of training-run variability. The result remains conditional on the selected synthetic WIC probability bank and static trust topology. The checkpoint, cache and input hashes, exact seed sets, all replicate outcomes, and summary are in \texttt{results/HIM/paired\_selection\_ciao\_wic\_train30\_seed42\_2026-10-11.json}.

I repeated the same configuration for training seeds 43 and 44, each with 30 instances, and independently evaluated each saved checkpoint for another 1,000 matched-seed replicates per selector. Best epochs were 19 and 26, respectively. Across the three training runs, ASW averaged 735.034 activated users (10.0456\% coverage) and MD averaged 734.662 (10.0405\%). The mean per-training-seed paired difference was $+0.371$ users, with sample standard deviation $0.531$ and a 95\% $t$ interval of $[-0.949,1.691]$ with only two degrees of freedom. ASW's per-run difference was positive for seeds 42 and 43 and slightly negative for seed 44. This initial result is consistent with no clear selection-method difference; the very small number of training runs leaves substantial uncertainty. Aggregate results and the three full replicate artifacts are recorded in \texttt{results/HIM/ciao\_wic\_three\_training\_seeds\_2026-10-11.json}.

\section{Synthesis and Project Direction}
The literature suggests three linked requirements: estimate influence from appropriate information, select complementary seeds, and evaluate the entire set transparently. Hyperbolic learning addresses the first through geometric structure; ASW is one attempt at the second. DeepSN and DeepIM offer alternative representation and optimization choices. Bandit and temporal studies identify regimes that a static reproduction does not cover. These are complementary questions rather than a progression in which the newest publication makes earlier methods unnecessary.

The next stage is to evaluate HIM on the audited Ciao trust graph, with DeepSN on the same user graph where its input and diffusion protocol can be aligned. The baseline Ciao release contains 7,317 users with trust relations and 85,205 edges after undirected projection, but its user ratings are not diffusion cascades and neither rating events nor trust links have usable timestamps in this release \cite{ciao}. HIM's graph adapter and fixed synthetic WIC bank now execute end-to-end, and three independent 30-instance runs are complete. The across-run ASW--MD interval includes zero and remains imprecise at three training seeds. Next add degree and random controls and coordinate the same user-level evaluation contract with DeepSN; consider a larger training-seed count and 1\%/10\% budgets only after the team protocol is aligned.

For a defensible result, freeze and hash the training cascades, edge-probability bank, node mapping, and evaluation seed bank. Report coverage with uncertainty across training seeds, while keeping Monte Carlo uncertainty for each fixed seed set separate. Record training time, seed-selection time, and evaluation time independently. For HIM specifically, verify that the selected nodes are unique and that the seed count equals the declared budget; save the exact checkpoint used for selection, the seed list, and the independent evaluation report. A useful research question is whether ASW's gain over MD is consistent on Ciao's trust graph, or whether the heuristic's neighbor-only overlap correction misses nonlocal audience overlap. Negative or mixed results would still identify a meaningful limitation of the selection mechanism.

HIM is a suitable reproduction target because its geometric interpretation and modular selection strategy generate specific, testable questions. Present experiments justify beginning the related-work report now; a literature review does not require completed benchmark comparisons. Stronger experimental conclusions should be added only when matched baselines and repeated trials support them.

\begin{thebibliography}{99}
\bibitem{kempe} D. Kempe, J. Kleinberg, and E. Tardos, ``Maximizing the spread of influence through a social network,'' in \emph{KDD}, pp. 137--146, 2003. \href{https://doi.org/10.1145/956755.956769}{doi: 10.1145/956755.956769}.
\bibitem{celf} J. Leskovec, A. Krause, C. Guestrin, C. Faloutsos, J. VanBriesen, and N. Glance, ``Cost-effective outbreak detection in networks,'' in \emph{KDD}, pp. 420--429, 2007. \href{https://doi.org/10.1145/1281192.1281239}{doi: 10.1145/1281192.1281239}.
\bibitem{discount} W. Chen, Y. Wang, and S. Yang, ``Efficient influence maximization in social networks,'' in \emph{KDD}, pp. 199--208, 2009. \href{https://doi.org/10.1145/1557019.1557047}{doi: 10.1145/1557019.1557047}.
\bibitem{deepim23} C. Ling, J. Jiang, J. Wang, M. T. Thai, R. Xue, J. Song, M. Qiu, and L. Zhao, ``Deep graph representation learning and optimization for influence maximization,'' in \emph{ICML}, PMLR 202, pp. 21350--21361, 2023. \href{https://proceedings.mlr.press/v202/ling23b.html}{Publisher paper}.
\bibitem{deepim24} T. Chowdhury, C. Ling, J. Jiang, J. Wang, M. T. Thai, and L. Zhao, ``Deep graph representation learning for influence maximization with accelerated inference,'' \emph{Neural Networks}, vol. 180, art. 106649, 2024. \href{https://doi.org/10.1016/j.neunet.2024.106649}{doi: 10.1016/j.neunet.2024.106649}.
\bibitem{deepsn} A. Hevapathige, Q. Wang, and A. N. Zehmakan, ``DeepSN: A sheaf neural framework for influence maximization,'' in \emph{AAAI}, vol. 39, no. 16, pp. 17177--17185, 2025. \href{https://doi.org/10.1609/aaai.v39i16.33888}{doi: 10.1609/aaai.v39i16.33888}.
\bibitem{bandit} Y. Feng, V. Y. F. Tan, and B. Cautis, ``Influence maximization via graph neural bandits,'' in \emph{KDD}, pp. 771--781, 2024. \href{https://doi.org/10.1145/3637528.3671983}{doi: 10.1145/3637528.3671983}.
\bibitem{poincare} M. Nickel and D. Kiela, ``Poincar\'e embeddings for learning hierarchical representations,'' in \emph{NeurIPS}, vol. 30, 2017. \href{https://proceedings.neurips.cc/paper/2017/hash/59dfa2df42d9e3d41f5b02bfc32229dd-Abstract.html}{Publisher paper}.
\bibitem{lorentz} M. Nickel and D. Kiela, ``Learning continuous hierarchies in the Lorentz model of hyperbolic geometry,'' in \emph{ICML}, PMLR 80, pp. 3779--3788, 2018. \href{https://proceedings.mlr.press/v80/nickel18a.html}{Publisher paper}.
\bibitem{him} H. Qiao, S. Feng, M. Zhou, X. Li, Y. Ye, F. Li, S. Shang, and Y.-S. Ong, ``Influence strength estimation in hyperbolic space for social influence maximization,'' \emph{IEEE Transactions on Knowledge and Data Engineering}, 2026. \href{https://doi.org/10.1109/TKDE.2026.3683562}{doi: 10.1109/TKDE.2026.3683562}; \href{https://arxiv.org/html/2502.13571v2}{full arXiv manuscript, v2}.
\bibitem{thgnets} Y. Liu, P. Zhang, W. Song, Y. Zheng, D. Li, L. Shi, and J. Gong, ``THGNets: Constrained temporal hypergraphs and graph neural networks in hyperbolic space for information diffusion prediction,'' in \emph{AAAI}, vol. 39, no. 12, pp. 12220--12228, 2025. \href{https://doi.org/10.1609/aaai.v39i12.33331}{doi: 10.1609/aaai.v39i12.33331}.
\bibitem{temporal} E. Yanchenko, T. Murata, and P. Holme, ``Influence maximization on temporal networks: A review,'' \emph{Applied Network Science}, vol. 9, art. 16, 2024. \href{https://doi.org/10.1007/s41109-024-00625-3}{doi: 10.1007/s41109-024-00625-3}.
\bibitem{ciao} J. Tang, H. Gao, H. Liu, and A. Das Sarma, ``eTrust: Understanding trust evolution in an online world,'' in \emph{KDD}, pp. 253--261, 2012; dataset page, \href{https://www.cse.msu.edu/~tangjili/datasetcode/truststudy.htm}{Ciao and Epinions data}.
\end{thebibliography}
\end{document}
